\subsection{Original PL data produce the full actual normal-model metric}

\begin{theorem}[Whole genuine normal ratio metric of an original PL function]
\label{local:actual-pl-whole-normal-ratio-metric}
\lean{EnvelopeProperty.ProjSubvariety.IsPL.exists_full_normalizedRatioMetric_models}
\leanok
\uses{local:actual-pl-all-product-point-monomial-models, local:actual-normalized-full-point-ratio-weight}
For every ORIGINAL PL function $\varphi$ on an original projective
subvariety $X$ over any field, there exist $r>0$, an integer $C$,
and two actual ORIGINAL product ideals $I,J$ such that their true logs
are finite at EVERY genuine Gauss-base product point. Consequently, on
the SAME actual original normalized simultaneous model built from these
ideals, for EVERY genuine relative valued point $y$ the TRUE negative-log
norm of the SAME canonical ratio section satisfies
$r\varphi(\operatorname{res}y)=\operatorname{weight}_{I/J}(y)+C$.
The whole model, invertible line, actual section, center norm and every
nonvanishing assertion are derived from original PL data. No additional
model, metric-identity, nefness or positivity certificate is assumed.
\end{theorem}
\begin{proof}
Use the actual original PL monomial model theorem, which produces genuine
original ideals with retained finite monomial generators and proves
finite logarithms and the exact PL identity at every original product
point. Take their actual native geometric restrictions to the original
variety product. Evaluate the derived all-point finiteness at original
Gauss lifts to obtain the inputs of the SAME existing normal simultaneous
model construction. For any original point $y$ of this actual model, the
all-point finiteness applied to its entire projected product point gives
nonvanishing of the actual first factor. The true full-point canonical
ratio-weight theorem identifies the actual section weight with the two
original ideal logarithms there. Substitute the already derived original
PL identity at that projected point. This proves the displayed equality
on the whole actual point space, without extrapolating from Gauss lifts.
\end{proof}
